\scanpage{22}
\mspage{11}
un syst. transitif d'iso entre les $\pi_t$ $(t\in\tau)$, permettant de les identifier tous à un $=T$. Ainsi $\langle\tau,\pi\rangle$ devient un $T$-groupoïde, \unsure{i.e.|c.àd.} un groupoïde (pas néc. connexe) où les grpes fondamentaux des comp. conn. sont comm. et tous identifiés à $T$.
De plus, $\mathcal E$ opère sur ce $T$-groupoïde (en opérant également sur $T$) via ses opérations sur le système $(\tau,\pi)$ ($\pi$-\unsure{espace|ens} à gauche) en opérant sur $\tau$ par l'op. donnée et sur $\pi$ via autom int. de $\mathcal E$, enfin sur $T$ via $\mathcal G$ via $\chi\ldots.$

[Si $\mathcal E$ opère sur $\tau$, \unsure{restreignant} cette op. à $\pi$ on trouve un groupoïde $\langle\tau,\pi\rangle$ sur lequel $\mathcal E$ opère. Supposant ce groupoïde abélien, on a donc une famille de groupes abéliens \inserted{$(T_j)_{j\in J}$} indexée par $J=\tau/\pi$, sur laquelle opère $\mathcal E$ via $\mathcal E/\pi=\mathcal G$. Ceci dit, on suppose donné une trivialisation de cette famille $(T_j)$ par $T\simeq T_j$, compatible avec op. de
\mspage{21}
